\subsection{Simulación de Eventos con Colas (escalera doble)}

\graybold{Preguntas guía:}

\begin{itemize}
\item ¿Hay un recurso compartido con estado (dirección, disponibilidad) que cambia según el orden de llegada de eventos con timestamps?
\item ¿Los eventos ya vienen ordenados (o se pueden ordenar) y conviene procesarlos con colas separadas por tipo?
\item ¿El estado siguiente depende de comparar el próximo evento contra un tiempo límite (fin de uso actual del recurso)?
\end{itemize}

\graybold{Utilidad:} Simular un sistema con estado (aquí: una dirección de movimiento y un instante de fin de uso) que evoluciona a medida que llegan eventos de distintos tipos, usando una cola por tipo de evento en vez de simular minuto a minuto.

\graybold{Complejidad:} \bigO{n} - cada persona se encola y desencola una sola vez.

\graybold{Fundamento teórico:} Se mantienen dos colas (una por dirección deseada) y un estado (dirección activa, instante en que el recurso queda libre). En cada paso se mira el evento más próximo entre ambas colas: si coincide con la dirección activa, extiende el fin de uso; si es la dirección contraria, esa persona debe esperar (se cuenta pero no se procesa aún); si ya no hay eventos antes de que el recurso quede libre, se decide si invertir la dirección (si había gente esperando) o terminar ese ciclo de uso.

\graybold{Ejercicio aplicado}

N personas llegan a una escalera doble (10 segundos por cruce) en tiempos distintos, cada una queriendo ir en una de dos direcciones. Determinar el instante en que la última persona termina de usarla.

\graybold{I/O:} n ≤ \ten{4} → lectura total + tokenización (patrón 2); se usan dos deques (collections.deque) en vez de queue, ya que en Python deque tiene \bigO{1} real para popleft().

\graybold{Código solución}

\begin{lstlisting}[language=Python]
import sys
from collections import deque
 
def solve():
    data = sys.stdin.buffer.read().split()
    idx = 0
    n = int(data[idx]); idx += 1
    left = deque()    # personas que quieren ir en direccion 0
    right = deque()   # personas que quieren ir en direccion 1
    for _ in range(n):
        t = int(data[idx]); d = int(data[idx + 1]); idx += 2
        (left if d == 0 else right).append(t)
 
    end = 0          # instante en que el recurso queda libre
    ans = 0
    waiting = 0       # cuantos esperan en direccion contraria a la activa
    direction = -1    # -1 = detenida
 
    while left or right or direction != -1:
        # mira el proximo evento disponible entre ambas colas
        next_time, who = None, -1
        if left and (next_time is None or left[0] < next_time):
            next_time, who = left[0], 0
        if right and (next_time is None or right[0] < next_time):
            next_time, who = right[0], 1
 
        if direction == -1:
            # arranca detenida: la toma quien llegue primero
            direction = who
            end = next_time + 10
            (left if who == 0 else right).popleft()
        elif who != -1 and next_time < end:
            if who == direction:
                end = max(end, next_time + 10)   # se sube mientras sigue en marcha
            else:
                waiting += 1                     # direccion contraria: debe esperar
            (left if who == 0 else right).popleft()
        else:
            # no hay mas eventos antes de que el recurso quede libre
            if waiting > 0:
                direction = 1 - direction   # invierte sentido para atender a los que esperan
                end += 10
                waiting = 0
            else:
                ans = end
                direction = -1               # se detiene
    print(ans)
solve()
\end{lstlisting}
